%! Setting Up a Test for the Difference of Two Population Means — Unit 7, Lesson 7.8 — Group Activity
\documentclass[10pt]{article}
\usepackage{apstats-article}
\usepackage{apstats-boxes}
\newcommand{\TallMath}[1]{$\displaystyle #1\rule[-1.4em]{0pt}{3.2em}$}

\begin{document}

\pageheader{Unit 7, Lesson 7.8}{Group Activity}
\namepartnerperiod

\vspace{0.12in}

% ── Tier R ─────────────────────────────────────────────────────────────────────
\begin{scenariobox}[Tier R --- Remediate: Match the Hypotheses]{sky}
\textbf{Directions:} Match each research scenario (1--3) to the correct pair of hypotheses
(A--C). Write the letter in the blank.

\medskip
\begin{minipage}[t]{0.48\linewidth}
\textbf{Scenarios}
\begin{enumerate}
  \item A researcher asks: ``Do students who listen to music while studying score
        \emph{differently} on average than those who study in silence?'' \quad \blank{1cm}
  \item A company claims its new battery lasts \emph{longer} on average than the
        standard brand. \quad \blank{1cm}
  \item A public health analyst believes mean daily step count is \emph{lower} for
        office workers than for retail workers. \quad \blank{1cm}
\end{enumerate}
\end{minipage}
\hfill
\begin{minipage}[t]{0.48\linewidth}
\textbf{Hypothesis Pairs}
\begin{itemize}
  \item[\textbf{A.}] $H_0: \mu_1 - \mu_2 = 0$; \; $H_a: \mu_1 - \mu_2 > 0$
  \item[\textbf{B.}] $H_0: \mu_1 - \mu_2 = 0$; \; $H_a: \mu_1 - \mu_2 \ne 0$
  \item[\textbf{C.}] $H_0: \mu_1 - \mu_2 = 0$; \; $H_a: \mu_1 - \mu_2 < 0$
\end{itemize}
\end{minipage}

\medskip
For Scenario 1, a conditions checklist is provided. Check each box if the condition is met
and write how you would verify it.

\smallskip
\begin{tabularx}{\linewidth}{@{} l X @{}}
$\square$ Random (Group 1): & \writeline \\
$\square$ Random (Group 2): & \writeline \\
$\square$ Normal (both groups are large, $n \ge 30$, OR distributions are checked): & \writeline \\
\end{tabularx}
\end{scenariobox}

% ── Tier A ─────────────────────────────────────────────────────────────────────
\begin{scenariobox}[Tier A --- Approaching Proficiency: Full Setup]{sky}
\textbf{Directions:} For the scenario below, (a) name the inference method, (b) state $H_0$
and $H_a$ in terms of $\mu_1 - \mu_2$, and (c) verify both conditions.

\medskip
\textbf{Scenario:} An education researcher randomly selects 32 students from classrooms using
project-based learning (Group 1) and 29 students from traditional lecture classrooms (Group 2)
at the same school. She records end-of-year standardized test scores. She believes project-based
learning leads to \emph{higher} mean scores. Histograms of both groups appear approximately
symmetric with no outliers.

\begin{enumerate}[label=\alph*.]
  \item Inference method: \blank{5.5cm}
  \item Define $\mu_1$ and $\mu_2$:

        $\mu_1 = $ \blank{6cm} \qquad $\mu_2 = $ \blank{6cm}

  \item $H_0$: \blank{4.5cm} \qquad $H_a$: \blank{4.5cm}

        Direction of $H_a$: \blank{2cm} \quad Justify: \writeline

  \item Conditions:
        \begin{itemize}
          \item Random (both groups): \writeline
          \item 10\% rule: \writeline
          \item Normal ($n_1 = 32$, $n_2 = 29$): \writeline
        \end{itemize}
\end{enumerate}
\end{scenariobox}

% ── Tier E ─────────────────────────────────────────────────────────────────────
\begin{scenariobox}[Tier E --- Extension: Design \& Compare]{sky}
\textbf{Directions:} Design a study and analyze the choice of alternative hypothesis.

\medskip
\textbf{Part 1 --- Study Design.} A health researcher wants to compare the mean number of
hours of sleep per night between college athletes and non-athletes at a large university.
Describe a study design (who would be sampled, how many, and why it produces independent groups).
\writelines{2}

\textbf{Part 2 --- Write the hypotheses.} Define $\mu_1$ (athletes) and $\mu_2$ (non-athletes)
and write $H_0$ and $H_a$ if the researcher believes athletes sleep \emph{less} on average.

$H_0$: \blank{4.5cm} \qquad $H_a$: \blank{4.5cm}

\textbf{Part 3 --- One-sided vs.\ two-sided.} If the researcher instead used $H_a: \mu_1 - \mu_2 \ne 0$,
how would this change the conclusion she could draw? What are the trade-offs between one-sided
and two-sided alternatives?
\writelines{3}
\end{scenariobox}

\begin{reflectionbox}
In Tier A, $n_2 = 29 < 30$. How does this change the Normal condition check compared to if
both groups had $n \ge 30$? Why must we examine \emph{both} distributions, not just the
smaller one?
\writelines{2}
\end{reflectionbox}

\end{document}
